\documentclass[11pt]{article}
\usepackage{lmodern}
\usepackage[T1]{fontenc}
\usepackage{amsmath,amssymb,amsthm}
\usepackage[margin=1in]{geometry}
\usepackage[hidelinks]{hyperref}

\newtheorem{theorem}{Theorem}
\newtheorem{lemma}[theorem]{Lemma}

\newcommand{\Q}{\mathrm Q}
\newcommand{\F}{\mathrm F}
\newcommand{\Z}{\mathrm Z}
\newcommand{\Ocal}{\mathcal O}
\newcommand{\Tr}{\operatorname{Tr}}
\newcommand{\Nm}{\operatorname{N}}
\newcommand{\Frob}{\operatorname{Frob}}

\title{An Explicit Quartic Ray Character
and a Frobenius Sign}
\author{\small\href{mailto:admin@cosmosgeek.org}{admin@cosmosgeek.org}}
\date{Working draft for specialist review}

\begin{document}
\maketitle

\begin{abstract}
We identify a quadratic character of a totally real
quartic field with an explicit mod-$4$ coefficient
functional. At primes having a unique degree-one
factor, this character equals a Legendre symbol
arising from a binary quartic covering of
$y^2=x^3-28x+1$.
The character also has a four-dimensional induced
Artin representation. The result is independent
of any analytic central-value conjecture. Its Galois closure has group
$W(D_4)$. The two signs on eligible primes are equidistributed
even after fixing a compatible cyclotomic residue class.
\end{abstract}

\section{The quartic field}

Put
\[
P(U)=U^4+3U^3-6U^2-4U+1,
\qquad L=\Q(u),\quad P(u)=0.
\]
Direct calculation gives
\[
\operatorname{disc}(P)=87781=41\cdot2141.
\]
Once irreducibility is verified below, squarefreeness gives
\[
\Ocal_L=\Z[u].
\]

The polynomial has four distinct real roots,
in the intervals
\[
(-5,-4),\quad(-1,-2/3),\quad
(0,1),\quad(1,2).
\]
Thus $L$ is totally real.

Modulo $2$,
\[
P(U)\equiv U^4+U^3+1,
\]
which is irreducible over $\F_2$ (there is no linear or irreducible
quadratic factor). This also establishes irreducibility over $\Q$.
Therefore $2$ is inert in $L$ and
\[
\Ocal_L/2\Ocal_L\simeq\F_{16}.
\]

\begin{lemma}\label{lem:classnumber}
The ordinary ideal class group of $L$ is trivial.
\end{lemma}

\begin{proof}
The Minkowski bound is
\[
\frac{4!}{4^4}\sqrt{87781}<28.
\]
It suffices to check prime ideals of norm
at most $27$.

The prime $2$ is inert, so its prime ideal
is $(2)$, of norm $16$.

The degree-one prime ideals above
$3,5,11,19,23$ are principal, as witnessed
by
\[
\begin{array}{c|c|r}
p&\alpha&\Nm_{L/\Q}(\alpha)\\ \hline
3&1+u&-3\\
5&1-u&-5\\
11&u^2-3&-11\\
19&1+2u&19\\
23&2+u&-23
\end{array}
\]
The degree-three prime above $3$ is also
principal, being the complementary factor
of $(3)$.

Modulo $7$, the quartic factors into two
irreducible quadratics, whose associated
prime ideals have norm $49$.
It is irreducible modulo $13$ and $17$,
so no prime ideal of norm $13$ or $17$
occurs.

These are all prime ideals below the
Minkowski bound. Hence every ideal class
is principal.
\end{proof}

\section{The governing quadratic extension}

Define
\[
d=u(3u+2)\in L^\times.
\]
Its norm is
\[
\Nm_{L/\Q}(d)=25.
\]

All four real embeddings of $d$ are positive.
Indeed, every real root of $P$ lies either
below $-2/3$ or above $0$, and
\[
x(3x+2)>0
\]
on those intervals.

Modulo $5$ the polynomial has exactly
one degree-one factor, at $u=1$.
The element $d$ has valuation two at
that prime, and valuation zero at every
other prime above $5$.

Because $\Nm(d)=25$, its valuation
at every odd prime ideal of $L$ is even.
Consequently the quadratic extension
\[
M=L(\sqrt d)
\]
is unramified at all odd finite primes.

It is also unramified at infinity,
since $d$ is totally positive.
Its finite conductor therefore divides
$4\Ocal_L$.

Let
\[
\chi_d:G_L\longrightarrow\{\pm1\}
\]
be the associated quadratic character.

\section{The exact dyadic character}

Write
\[
\mathrm F=\Ocal_L/2\Ocal_L=\F_{16},
\]
and denote reduction modulo $2$
by a bar.

Because $u$ is a $2$-adic unit,
\[
\frac d{u^2}
=3+\frac2u
=1+2(1+u^{-1}).
\]

Set
\[
\beta=1+\bar u^{-1}\in\mathrm F.
\]
The identity
\[
\bar u^4+\bar u^3+1=0
\]
implies
\[
\bar u^{-1}=\bar u^3+\bar u^2,
\]
so
\[
\beta=1+\bar u^2+\bar u^3.
\]

Let $\alpha\in\Ocal_L$ be prime to $2$.
Since $\mathrm F^\times$ has order $15$,
\[
\alpha^{15}\equiv1\pmod2.
\]
There is therefore a well-defined
\[
c(\alpha)=
\frac{\alpha^{15}-1}{2}\pmod2
\quad\in\mathrm F.
\]

Write
\[
c(\alpha)=c_0+c_1\bar u+c_2\bar u^2
+c_3\bar u^3,
\qquad c_i\in\F_2,
\]
and define
\[
e(\alpha)=c_1\in\F_2.
\]

\begin{lemma}\label{lem:dyadic}
For every $\alpha\in\Ocal_L$ prime to $2$,
\[
(d,\alpha)_2=(-1)^{e(\alpha)},
\]
where the Hilbert symbol is taken in
the unique completion of $L$ above $2$.
\end{lemma}

\begin{proof}
The local field $L_2/\Q_2$ is unramified
of degree four.

Its odd-order Teichmuller units have trivial
quadratic Hilbert symbol against $d$.
Writing the principal-unit component
of $\alpha$ as $1+2a$, we have
\[
\overline a=c(\alpha),
\]
because raising to the fifteenth power
kills the Teichmuller component and
preserves the first principal-unit layer
modulo $4$.

Here is the local norm test. For the unit $b=1+u^{-1}$,
$\sqrt{1+2b}-1$ is Eisenstein over $L_2$.
Writing a unit in its quadratic integer ring as
$x+y(\sqrt{1+2b}-1)$ gives norm $x^2-2xy-2b y^2$.
A norm congruent to $1\bmod 2$ requires $x\equiv1\bmod2$.
The congruence for $1+2a\bmod4$ becomes
$\bar a=\bar y+\bar b\bar y^2$, soluble exactly when
$\Tr(\bar b\bar a)=0$, by the trace-zero image of $z\mapsto z+z^2$.
Conversely a solution lifts by Hensel's lemma to
$-y-by^2=a$, giving the required norm with $x=1$.
Thus
\[
(d,\alpha)_2
=(-1)^{\Tr_{\mathrm F/\F_2}
(\beta c(\alpha))}.
\]

A direct trace calculation in
\[
\mathrm F=\F_2[\bar u]/
(\bar u^4+\bar u^3+1)
\]
gives
\[
\begin{array}{c|cccc}
x&1&\bar u&\bar u^2&\bar u^3\\ \hline
\Tr(\beta x)&0&1&0&0.
\end{array}
\]
Therefore
\[
\Tr(\beta c(\alpha))=c_1=e(\alpha),
\]
which proves the assertion.
\end{proof}

The character is nontrivial on the
first principal-unit layer, so its
local conductor at $2$ is exactly $4$.

\paragraph{Uniqueness.}
The finite ray group modulo $4\Ocal_L$ is
$(\Ocal_L/4)^\times/\operatorname{im}\Ocal_L^\times$ since the
ordinary class number is one. Its numerator has order $240$.
The units $-1,-u,-1-u+u^2,-1+2u^2$ have norms $1,1,-1,-1$.
Adjoining their residues successively generates groups of orders
$2,30,60,120$. These small computations can be reproduced by starting
with $T=\{1\}$ and replacing it for each unit $g$ by
$\{t g^i:t\in T,\ 0\le i<240\}$, reducing coefficients modulo
four and powers of $u$ by $P$. The nontrivial dyadic character is
trivial on global units by Hilbert reciprocity. Hence the finite
ray group has order two and $M/L$ is its quadratic extension.

\section{The Frobenius-sign theorem}

Consider
\[
G_1(X,Z)=
-16X^4-24X^3Z+24X^2Z^2
+8XZ^3-Z^4.
\]
The defining polynomials satisfy
\[
P(2X)=-G_1(X,1).
\]

Let $q$ be an odd rational prime such that
\[
q\nmid 5\cdot41\cdot2141
\]
and $P$ has factorization type $(1)(3)$
modulo $q$.

Let $u_q\in\F_q$ be its unique root,
and put
\[
r_q=u_q/2.
\]
Then $r_q$ is the unique root of
$G_1(X,1)$ modulo $q$.

The corresponding degree-one prime ideal is
\[
\mathfrak q=(q,u-u_q).
\]

By Lemma~\ref{lem:classnumber}, choose
$\alpha_q\in\Ocal_L$ such that
\[
(\alpha_q)=\mathfrak q.
\]
Necessarily
\[
\Nm(\alpha_q)=\pm q.
\]

Define the Fisher-type quadratic sign
\[
\tau_q=
\left(\frac{r_q(3r_q+1)}q\right).
\]
The Legendre argument is nonzero for
the primes under consideration.

\begin{theorem}\label{thm:governing}
For every such prime $q$,
\[
\boxed{
\tau_q=(-1)^{e(\alpha_q)}.
}
\]
The exponent is independent of the
choice of generator $\alpha_q$.
\end{theorem}

\begin{proof}
Since $u_q=2r_q$,
\[
u_q(3u_q+2)
=4r_q(3r_q+1).
\]
Hence
\[
\tau_q=
\left(\frac{d\bmod\mathfrak q}{q}\right).
\]

Because $q$ is unramified in $M/L$,
the local norm-residue symbol at
$\mathfrak q$ is
\[
(d,\alpha_q)_{\mathfrak q}=\tau_q.
\]

At every other odd finite place,
$\alpha_q$ is a unit and $d$ has even
valuation. The corresponding local
Hilbert symbol is therefore $1$.

At every real place, $d>0$, so the
Hilbert symbol is again $1$.

Hilbert reciprocity now gives
\[
1=\prod_v(d,\alpha_q)_v
=(d,\alpha_q)_2\,\tau_q.
\]
By Lemma~\ref{lem:dyadic},
\[
(d,\alpha_q)_2=(-1)^{e(\alpha_q)}.
\]
Since both values are signs, the
claimed identity follows.

For generator independence, suppose
$\alpha'_q=\varepsilon\alpha_q$
with $\varepsilon\in\Ocal_L^\times$.

Hilbert reciprocity applied to
$(d,\varepsilon)$ gives
\[
(d,\varepsilon)_2=1,
\]
because all its other local symbols
are trivial.

Lemma~\ref{lem:dyadic} implies
$e(\varepsilon)=0$.
The map $e$ is additive on odd units:
\[
e(\alpha\beta)
=e(\alpha)+e(\beta)\pmod2,
\]
as follows by multiplying their
fifteenth powers modulo $4$.

Therefore
\[
e(\alpha'_q)=e(\alpha_q).
\]
\end{proof}

\section{The induced Artin trace}

Let
\[
\rho_d=
\operatorname{Ind}_{G_L}^{G_\Q}\chi_d.
\]
This is a four-dimensional Artin
representation of $G_\Q$.

At a prime $q$ of quartic Frobenius
cycle type $(1)(3)$, Frobenius permutes
the four induced cosets with exactly
one fixed point.

The three-cycle contributes zero to
the trace. The fixed coset contributes
the local character at the unique
degree-one prime.

Consequently
\begin{equation}\label{eq:artin-trace}
\boxed{
\operatorname{Tr}
\rho_d(\Frob_q)
=
\chi_d(\Frob_{\mathfrak q})
=\tau_q=(-1)^{e(\alpha_q)}.
}
\end{equation}

Thus the quartic-root sign is a
conjugacy-invariant Artin trace.

\section{Two exact controls and density}

At $q=23$, the unique quartic root is
$u_{23}=-2$, with generator
\[
\alpha_{23}=-2-u,\qquad
\Nm(\alpha_{23})=-23.
\]
Direct calculation modulo $4$ gives
\[
\alpha_{23}^{15}
\equiv3+2u^2+2u^3\pmod4.
\]
Hence $e(\alpha_{23})=0$ and
\[
\tau_{23}=+1.
\]

At $q=31$, the unique quartic root is
$u_{31}=6$, with generator
\[
\alpha_{31}=-3-2u+3u^2,\qquad
\Nm(\alpha_{31})=31.
\]
Here
\[
\alpha_{31}^{15}
\equiv1+2u+2u^2+2u^3\pmod4.
\]
Thus $e(\alpha_{31})=1$ and
\[
\tau_{31}=-1.
\]

Both primes satisfy
\[
q\equiv7\pmod8,\qquad
\left(\frac{-q}{41}\right)
=\left(\frac{-q}{2141}\right)=1.
\]

Let $\widetilde M$ be the Galois
closure of $M/\Q$, and put
\[
H=\widetilde M\,
\Q(\zeta_{8\cdot41\cdot2141}).
\]

The extension $H/\Q$ is finite Galois.
Both $23$ and $31$ are unramified in $H$.

The indicated splitting conditions
and the sign in
\eqref{eq:artin-trace}
are constant on Frobenius conjugacy
classes in $\operatorname{Gal}(H/\Q)$.

Chebotarev therefore gives
positive-density sets of primes
satisfying the displayed local
conditions and having, respectively,
\[
\tau_q=+1
\quad\text{and}\quad
\tau_q=-1.
\]

\section{The full sign group and the two densities}

Let $V_0=\{e\in\F_2^4:\sum e_i=0\}$. The Galois group of $P$ is
$S_4$: its Frobenius permutations at $2$ and $23$ contain respectively
a four-cycle and a three-cycle. The order of the group is divisible
by $12$ and the unique index-two subgroup of $S_4$ contains no
four-cycle. Since $\prod_i d(u_i)=25$ we may choose the four square
roots with product $5$. Thus the Galois group of $\widetilde M$ embeds
in the even signed permutation group $V_0\rtimes S_4$. Modulo $277$
the roots $u_i$ are $5,71,229,246$, with signs of $d(u_i)$ equal to
$+,+,-,-$, respectively. This prime divides neither $2\cdot 5$ nor the
quartic discriminant. Its Frobenius over the identity permutation is
a two-coordinate flip. Conjugating supplies every such flip, and
they span $V_0$. The embedded group thus has order $8\cdot24=192$
and is the whole $W(D_4)$.

The commutators of $S_4$ with $V_0$ span $V_0$ (for instance
the difference of the pair vectors $e_i+e_j,e_i+e_k$ is $e_j+e_k$).
The maximal abelian subfield of $\widetilde M$ is therefore
$\Q(\sqrt{87781})$. It is contained in $\Q(\zeta_B)$ for
$B=8\cdot41\cdot2141$, since $87781$ is a positive fundamental
discriminant dividing $B$. The compositum has order $96\varphi(B)$.
Fix a unit residue class modulo $B$ acting trivially on this
discriminant field. Among the $96$ compatible lifts lie eight
quartic three-cycles with eight sign vectors apiece. For each of
these three-cycles exactly four sign vectors have each value at the
fixed root. Consequently the natural density among all primes of
each joint class (specified residue, quartic type $(1)(3)$ and
specified sign) is
\[
 \frac{32}{96\varphi(B)}=\frac1{1\,027\,200}.
\]
In particular this applies to the class of $31$ modulo $B$.

Let $z=(1,1,1,1)\in V_0$ be the central flip. Over characteristic
zero the restriction of any representation to $V_0$ decomposes
into characters, indexed by subsets of $\{1,2,3,4\}$ up to complement.
The $S_4$-orbits of these characters have sizes $1,3,4$.
Exactly the four singleton characters are negative on $z$, so any
representation on which $z$ acts nontrivially has dimension at
least four. The signed permutation representation realizes equality.

There is a complementary finite-field obstruction. Consider a
finite compositum $J$ of cyclotomic fields and fields cut out by
two-dimensional representations over finite fields of characteristic
two. The Sylow two-subgroups of each $\operatorname{GL}_2(\F_{2^f})$
are abelian (upper unitriangular), and the same holds for subgroups
of their product with an abelian group and for their quotients.
But $W(D_4)$ has a nonabelian Sylow two-subgroup, as a double
transposition does not commute with every pair flip. Thus
$\widetilde M$ is not contained in $J$. Every nontrivial normal
subgroup of $W(D_4)$ contains $z$: if its intersection with $V_0$
is nontrivial it contains $z$ or a pair flip and hence all pair
flips; if the intersection is trivial the subgroup centralizes
$V_0$, whose $S_4$-action is faithful, so the subgroup is trivial.
Enlarge $J$ by $\Q(\zeta_B)$. In the compositum one can therefore
change a lift of the Frobenius of $31$ by $(z,1)$ while fixing its
$J$-component. Chebotarev gives positive-density sets of primes
with opposite signs but the same Frobenius conjugacy class on $J$.

Finally the kernel of $W(D_4)\to S_4/V_4=S_3$ is
$V_0\rtimes V_4$, where $V_4$ is the normal Klein four group.
The commutators $v+\sigma v$ for $\sigma\in V_4$ and $v\in V_0$
belong to $\{0,z\}$. The action of $V_4$ on $V_0$ is faithful and
its fixed subspace is $\langle z\rangle$. Thus the center and the
derived subgroup both equal $\langle z\rangle$, and the quotient
by it is elementary abelian of order $16$. This gives an extraspecial
group of order $32$, of plus type as $V_0$ is elementary abelian of
order $8$.

\section{Relation to descent and scope}

For $q\equiv7\bmod 8$ with $(-q/41)=(-q/2141)=1$ and quartic type
$(1)(3)$, the separate binary-quartic calculation in the
\href{https://github.com/cosmosgeek0/BSD/blob/main/fisher-twist-rank-zero.pdf}{prime-twist manuscript}
identifies this sign with the Cassels--Tate pairing of its two
specified classes. The ray digit and Artin trace here describe the
arithmetic sign. The conjectural congruence
$I_q\equiv12e(\alpha_q)\bmod24$ for the corresponding normalized
central value remains open for the full family.

\begin{thebibliography}{9}
\raggedright
\bibitem{Fisher} T.~Fisher, \emph{On binary quartics and the Cassels--Tate pairing},
Res. Number Theory \textbf{8} (2022), article 74.
\href{https://doi.org/10.1007/s40993-022-00376-z}{Publisher's paper}.
\bibitem{Smith} A.~Smith, \emph{Governing fields and statistics for
4-Selmer groups and 8-class groups},
\href{https://arxiv.org/abs/1607.07860}{arXiv:1607.07860}.
\end{thebibliography}
\end{document}

